\section{动机悖论与冲动理性：你的奖励系统长在别处}

上一章讲了能量，这一章讲动机——为什么「重要」的事推不动，而「没意义」的事反而写得飞快。

\subsection{动机悖论：压力越大，越启动不了}

原视频讲过一个让很多人拍大腿的例子（BV1iKPrzNEjN）：公司让你写一本书，告诉你写好了能卖钱——你一直拖。但如果换个念头：「人类要完蛋了，这本书写不写无所谓，随便写写图个好玩」——你反而很快就写完了，质量还不错。

从逻辑上看这是个悖论：你都觉得人类要完蛋了，还写什么书？但从神经机制上看，一点都不悖论：

\term{「能卖钱、要交差、写不好会被评判」这些社会规范压力，会推高你的皮质醇。皮质醇一高，启动就锁死了。而「人类要完蛋」这个念头，功能上消除了社会规范压力，皮质醇降下来，你反而能动了。}

原视频为此给了一个极重要的心智工具：\term{一元论视角}。不存在一个坐在大脑里发号施令的「统一的自我」；信念不是指令，而是网络状态。一个念头（「人类要完蛋」）的价值不在于它是不是真理，而在于它产生了什么效果——系统只问效用，不问真理。只要它好用，就把它留着，不必要求自己逻辑自洽。

\begin{insight}
这对 ADHD 自救的意义非常大：\term{别执着于「我应该是自洽的、认真的、负责的」}。那些宏大叙事——「我要做一个靠谱的人」「这份工作必须一丝不苟地完成」——很可能正是锁死你启动的东西。该自我欺骗的时候就自我欺骗，只要它真的降低焦虑、释放执行力，从神经学的角度看，它就是真实。
\end{insight}

\subsection{冲动理性 vs 时间理性}

普通人的「理性」，是一种\term{时间理性}（BV1apQ2ByEN3）：假设干 A 还是干 B，区别只有那点切换成本，所以按部就班永远最优。

但 ADHD 的「理性」不是这样运行的。原视频用打游戏举例子：你砍一个血很厚的 BOSS，砍几刀就觉得无聊想走——时间理性会说「走了还得重新加载，浪费几十秒，不划算」；但\term{冲动理性}会说：「无聊了就走，因为现在有个更有价值的冲动冒出来了，冲动伴随新想法，走是对的，BOSS 以后再来打。」

对 ADHD 来说，\term{价值权重不在「这个程序花多少时间」，而在「当下的神经兴奋度，以及兴奋指向哪里」}。同样一份小组报告，冲动来的时候三十分钟写完；被硬安排着写，一天都写不完，还要边写边想「这破报告对我有啥用，写不好还要被嘲笑」。

所以：

\begin{checklist}
  \item \term{冲动不是敌人，是信号}。第一次冲动大多没用，但它是种子——二次冲动、三次冲动会长成深刻的问题。随手记录下来，别让它溜走（工具上 Obsidian 非常合适）。
  \item \term{把冲动引导成网，而不是捋成线}。你不可能把冲动排成一条直线，但可以引导成一张大致有用的网——在大方向内跳来跳去，而不是被任务死死锁死。
  \item \term{考试能专注的真相}：考试时你走不了，环境强制压制了其他冲动。但这个办法没法日常化——现代人有电脑有手机，把自己关进小黑屋也有一百种办法分心。
\end{checklist}

\subsection{兴趣引导：为什么「无敌」之后反而无聊}

单机游戏开挂变成无敌，很多人反而觉得没意思了——这几乎是所有人的体验，但原视频指出它和 ADHD 有深刻的联系（BV1Y2wizyEua）：\term{大脑在扫描「有没有值得做的事」，扫描结果是没有，于是陷入无聊}。游戏无敌后是这样；ADHD 面对「写了也没用」的任务时，也是这样。

区别在于：普通人能说服自己「这东西能卖钱、能交差」，硬着头皮驱动自己；ADHD 对这些外部驱动力不那么感冒，大脑的预测结果就是「这东西写了也没啥用」，于是干脆无法启动。

\begin{insight}
所以 ADHD 的正确姿势不是「逼自己启动一个写了也没用的任务」，而是\term{让自己尽可能待在一个可以跳跃的激发场里}——你定向驱动的能量非常有限，但你长期经营好一个场域，让有意义的问题在这个场域里冒出来的概率足够高，然后顺着它做。这叫\term{以广促深}，是贯穿这份讲义的总心法。
\end{insight}

\subsection{隐形回血：管理分心，而非消除分心}

最后补一个被低估的机制（BV16hxrzYEpT）：\term{全神贯注的效率，会被事后累抵消}。

让 ADHD 全神贯注地敲代码，任务期间确实更快，但敲完可能要躺三个小时——把效率全赔回去。反过来，允许一些「有价值的创新性分心」：敲代码时突然冒出一个想法，「这个定理如果和那个东西混合会怎么样」，花一分钟记下来或问一下 AI——这不只是回血，还是好想法的种子，甚至会让「启动这件事」的门槛降低（因为你知道中途可以合法地跑神）。

所以 ADHD 的目标不是\term{消除分心}，是\term{管理分心}——把分心的范围控制在有价值区间内：

\begin{checklist}
  \item \textbf{无新内容的冲动}（刷手机、开无关网页、想断连）：用身份认同压下去——「我是一个能按顺序行动的人」；
  \item \textbf{有点意思的想法}（冒出一个新的连接、新的猜想）：记录下来或问 AI，那是回血，也是种子。
\end{checklist}

\key{对 ADHD 谈「效率」，常常是复收益：没有快速反馈时，效率指标里全是「做了也没用」的评判。真正的奖励不是钱，是随机冒出来的好想法——那才是你的充电桩。}
